% 出席確認クイズ 第02回　情報と意思決定
%   attendance/attendance02.html から生成（解説は本ガイド用に加筆）。

\quizq{1}{データ・情報・知識の階層のうち、「データに文脈や意味を与えたもの」はどれでしょうか?}%
  {知恵}%
  {\correct{情報}}%
  {データ}%
  {雑音}%
  {正解は (b)。DIKW 階層では、データ（事実の記録）に文脈や意味を与えたものが情報、情報を体系化して活用できるようにしたものが知識、知識を状況に応じて使い分ける力が知恵です。たとえば「気温32度」はデータ、「例年より5度高い」は情報、「暑い日はアイスが売れるので仕入れを増やす」は知識にあたります。(d) の雑音はデータに混じる意味のない部分です。}

\quizq{2}{サイモンの意思決定プロセスの3つの段階に含まれ\underline{ない}ものはどれでしょうか?}%
  {選択活動(Choice)}%
  {\correct{宣伝活動}}%
  {設計活動(Design)}%
  {情報活動(Intelligence)}%
  {正解は (b)。サイモンは意思決定を情報活動（Intelligence）・設計活動（Design）・選択活動（Choice）の3段階で説明しました。後に再検討（Review）を加えることもあります。「宣伝活動」は含まれません。英語名と日本語名の対応を覚えておくと、文献を読むときに役立ちます。}

\quizq{3}{意思決定プロセスのうち、複数の代替案を作り出す段階はどれでしょうか?}%
  {選択活動(Choice)}%
  {\correct{設計活動(Design)}}%
  {情報活動(Intelligence)}%
  {評価活動}%
  {正解は (b)。設計活動は、問題に対する解決策の代替案を考え出す段階です。情報活動は問題や機会を発見する段階、選択活動は代替案の中から1つを選ぶ段階です。(d) の評価活動という名称は3段階にはありません。生成AIは代替案を数多く出せるため、設計活動との相性が特に良いと言えます。}

\quizq{4}{「限定合理性」の説明として適切なものはどれでしょうか?}%
  {\correct{人は情報処理能力の限界のもとで、満足できる水準の選択をする}}%
  {AIには合理性がまったくない}%
  {合理的な判断は経営には不要である}%
  {人はどんな状況でも常に最適な選択ができる}%
  {正解は (a)。限定合理性（bounded rationality）はサイモンの中心概念で、人は完全な情報と無限の処理能力を持たないため、最適解ではなく「満足できる解」を選ぶ（満足化）という考え方です。(d) は古典的な完全合理性の仮定です。(b)(c) は誤りで、AIにも学習データや設計に由来する限界があります。}

\quizq{5}{組織の階層と情報の関係として適切なものはどれでしょうか?}%
  {経営層は情報を必要としない}%
  {情報は階層に関係なく同じ形式で使われる}%
  {\correct{経営層ほど要約された非定型の情報を扱う傾向がある}}%
  {現場ほど要約された情報だけを使う}%
  {正解は (c)。アンソニーの3階層（戦略的計画・マネジメントコントロール・業務統制）に対応して、経営層ほど要約された非定型で外部志向・将来志向の情報を、現場ほど詳細で定型的な内部情報を扱う傾向があります。(a)(b)(d) はこの関係を否定するか逆にしています。}
